%! Author = lili
%! Date = 7/9/26

\section{Tutorium 29.06.2026}
\untit{Präsenzaufgaben 2.4: Erste Schritte in der Statistik \hfill 29.\ Juni -- 03.\ Juli 2026}

\txo{Bearbeiten Sie auf jeden Fall Aufgabe 2.4.4. Falls Ihnen Zeit bleibt, bearbeiten Sie
zusätzlich auch Aufgabe 2.4.5.}
\subsection{PÜ 2.4.4. (Lösung da, noch nicht eingetragen)}
% TODO Bilder Nils WhatsApp abtippen
Seien $X_1,\dots,X_n$ unabhängige, identisch verteilte Stichprobenvariable. Zeigen Sie:

\begin{auf}[Der Schätzer $t$, der gegeben ist durch $t(x) = \tfrac{1}{n}\sum_i x_i$ für alle
$x = (x_1,\dots,x_n)$, ist ein erwartungstreuer Schätzer für den Erwartungswert
der Stichprobenvariablen.]
    NOCH KEINE LÖSUNG % TODO NOCH KEINE LÖSUNG
\end{auf}

\begin{auf}[Wir setzen nun zusätzlich voraus, dass $\Var(X_1) < \infty$ gilt. Berechnen Sie
das Risiko des Schätzers $t$.]
    NOCH KEINE LÖSUNG % TODO NOCH KEINE LÖSUNG
\end{auf}

\begin{auf}[Wie verhält sich das Risiko von $t$ für $n \to \infty$?]
    NOCH KEINE LÖSUNG % TODO NOCH KEINE LÖSUNG
\end{auf}

\emph{Bitte beachten Sie:} Wir nehmen in dieser Aufgabe nicht an, dass es sich um
Bernoulli-Variable handelt.

\subsection{PÜ 2.4.5}
Ein kleiner Statistikkurs besteht aus sieben Studierenden. Am Ende des Kurses wird eine
Klausur geschrieben, bei der maximal $120$ Punkte erreicht werden können. Die
Studierenden erreichten die folgenden Punktzahlen:
\[
    \begin{array}{c|ccccccc}
        \text{Studierende(r)} & A & B & C & D & E & F & G \\ \hline
        \text{Punkte} & 97 & 80 & 15 & 25 & 55 & 68 & 80
    \end{array}
\]
Berechnen Sie:
\begin{auf}[das arithmetische Mittel der erreichten Punkte,]
    NOCH KEINE LÖSUNG % TODO NOCH KEINE LÖSUNG
\end{auf}

\begin{auf}[den Median der erreichten Punkte,]
    NOCH KEINE LÖSUNG % TODO NOCH KEINE LÖSUNG
\end{auf}

\begin{auf}[den Modalwert der erreichten Punkte,]
    NOCH KEINE LÖSUNG % TODO NOCH KEINE LÖSUNG
\end{auf}

\begin{auf}[die Spannweite der erreichten Punkte,]
    NOCH KEINE LÖSUNG % TODO NOCH KEINE LÖSUNG
\end{auf}

\begin{auf}[die mittlere absolute Abweichung vom Median der erreichten Punkte,]
    NOCH KEINE LÖSUNG % TODO NOCH KEINE LÖSUNG
\end{auf}

\begin{auf}[die Varianz und die Standardabweichung der erreichten Punkte.]
    NOCH KEINE LÖSUNG % TODO NOCH KEINE LÖSUNG
\end{auf}

Angenommen jede(r) der Studierenden hätte $3$ Punkte mehr erreicht. Wie groß wären dann

\begin{auf}[das arithmetische Mittel der erreichten Punkte?]
    NOCH KEINE LÖSUNG % TODO NOCH KEINE LÖSUNG
\end{auf}

\begin{auf}[der Median der erreichten Punkte?]
    NOCH KEINE LÖSUNG % TODO NOCH KEINE LÖSUNG
\end{auf}

\begin{auf}[die Varianz und die Standardabweichung der erreichten Punkte?]
    NOCH KEINE LÖSUNG % TODO NOCH KEINE LÖSUNG
\end{auf}

Berechnen Sie die Werte (g), (h) und (i) mit Hilfe der Ergebnisse in (a), (b) und (f).

% TODO von nils bildern noch die ÜB Besprechung abtippen